\scanpage{4}
-32-

In view of th.11, it is natural to define a homomorphism
\[
\rho_\infty:I_K\longrightarrow F_\infty^*=(R\otimes F)^*
\]
by putting $\rho_\infty(a)=\varepsilon(a)\psi_\infty(a_\infty^{-1})$, where $a_\infty$ is the infinite component of the idèle $a$. This homomorphism is obviously characterized by the fact that it is continuous, trivial on $K^*$, and coincides with $\varepsilon$ on the group $I_K^\infty$ of idèles whose infinite component is 1.

Let $\sigma:F\to C$ be a homomorphism. The composition
\[
\chi_\sigma:I_K\xrightarrow{\rho_\infty}(R\otimes F)^*\xrightarrow{1\otimes\sigma}C^*
\]
is continuous and trivial on $K^*$ ; that is, $\chi_\sigma$ is a ``Grössencharakter'' in the broad sense (having values in $C^*$ rather than in the unit circle); it is essentially the same as the Grössencharakter defined in [18], p.148. For each valuation $v$ of $K$, the restriction of $\chi_\sigma$ to $U_v(K)$ is $\sigma\circ\varphi_v$, and, since $\sigma$ is injective, it follows that \uline{the exponent at} $v$ \uline{of the conductor of} $\chi_\sigma$ \uline{is equal to the number} $n_v$ \uline{of th.}6. Hence :

THEOREM 12 - \uline{The conductor of the abelian variety} $A$ \uline{is the} $2d$-th \uline{power of the conductor of} $\chi_\sigma$ ; \uline{in particular}, \uline{the support of the conductor of} $\chi_\sigma$ \uline{is the set of valuations of} $K$ \uline{where} $A$ \uline{has bad reduction}.

For elliptic curves, the second statement was proved by Deuring [7].
\mspage{3}
